% §7.1 The Setup Revisited
%
% Cross-reference resolution status:
%
% DEFINED here:
%   \label{sec:setup-revisited}                 -- this section
%
% ALREADY DEFINED (chapter 1):
%   \Cref{sec:sampling-rates}                   -- chapter 1 §1.4
%   \Cref{sec:quantifying-noise}                -- chapter 1 §1.2
%
% ALREADY DEFINED (chapter 5):
%   \Cref{ch:approximate-sets}                  -- chapter 5
%   \Cref{def:bloom-filter}                     -- chapter 5 §5.2
%
% FORWARD REFS within chapter 7 (resolve in later Plan 7 tasks):
%   \Cref{sec:binomial-estimator}               -- §7.2 (Task 3)
%   \Cref{sec:confidence-intervals}             -- §7.3 (Task 4)
%   \Cref{sec:derived-measures}                 -- §7.4 (Task 5)
%   \Cref{sec:bloom-diagnostics}                -- §7.5 (Task 6)
%
% §7.1 opens chapter 7 by carrying chapter 1 §1.4's setup forward into
% Part II's notation and naming the four items §1.4 deferred. Short
% recap; no new theorems.

\section{The Setup Revisited}
\label{sec:setup-revisited}

\Cref{sec:sampling-rates} stated the problem and stopped just short of
the theory. An approximate set arrives with two unknown rates: a rate
$\fprate$ at which it reports $\top$ on inputs whose truth is $\bot$,
and a rate $\fnrate$ at which it reports $\bot$ on inputs whose truth
is $\top$. The classifier was treated as a black box: query it on $n$
labelled inputs, split them into the $n_0$ trials whose truth is
$\bot$ and the $n_1 = n - n_0$ trials whose truth is $\top$, count the
false positives and false negatives, and report the empirical
proportions $\hat{\fprate} = \mathrm{FP}/n_0$ and $\hat{\fnrate} =
\mathrm{FN}/n_1$. Part~II now writes that classifier as an approximate
set $\ASet{A} : \Universe \to \Bernoulli{\BoolType}$, a noisy
membership predicate over a universe $\Universe$, but the estimation
problem is unchanged. Chapter~7 closes the loop.

\subsection*{What chapter 1 left deferred}

\Cref{sec:sampling-rates} reported the binomial form of the
estimators, stated their unbiasedness and variance without proof, and
gave the Wald confidence interval with a warning attached. It promised
the full theory to chapter~7. Four items were named:

\smallskip
\noindent\emph{(i) Estimator properties.}
\Cref{sec:sampling-rates} claimed $\hat{\fprate} = \mathrm{FP}/n_0$ is
unbiased for $\fprate$ with variance $\fprate(1 - \fprate) / n_0$. It
did not prove either claim. It did not identify $\hat{\fprate}$ as the
maximum-likelihood estimator, and it did not show that the estimator's
variance equals the Cramer-Rao lower bound. All four facts (bias,
variance, MLE identity, efficiency) are stated and proved in
\Cref{sec:binomial-estimator}.

\smallskip
\noindent\emph{(ii) Proper confidence intervals.}
The Wald interval $\hat{\fprate} \pm z_{\alpha/2}
\sqrt{\hat{\fprate}(1 - \hat{\fprate}) / n_0}$ is asymptotically
correct but inadequate in the small-$n_0$ regime where the cost of a
single observation matters. Three standard alternatives exist (Wilson,
Clopper-Pearson, Jeffreys), each with a different motivation and a
different operating characteristic. \Cref{sec:confidence-intervals}
derives each and recommends a default.

\smallskip
\noindent\emph{(iii) Sawtooth under-coverage near the boundaries.}
The chapter~1 notebook visualized the Wald interval failing at small
$n_0$ when $\fprate$ is near $0$ or $1$: empirical coverage drops
below the nominal $1 - \alpha$ in a jagged, oscillating pattern that
worsens as the rate approaches a boundary. The pattern has a name
(sawtooth coverage) and a known cause (discrete binomial counts
viewed through a continuous Normal approximation). The same section
diagnoses the failure and shows which of the alternative intervals
suppress it.

\smallskip
\noindent\emph{(iv) Derived measures.}
The rates $\fprate$ and $\fnrate$ are intrinsic to a classifier;
precision, the positive predictive value (PPV), the negative
predictive value (NPV), and the F1 score are not. Each is a function
of $\fprate$, $\fnrate$, and the prevalence $\pi$ of $\top$-inputs in
the workload. \Cref{sec:derived-measures} writes each derived measure
in those three variables and uses the delta method to turn the
binomial CIs of \Cref{sec:confidence-intervals} into CIs for the
derived quantities.

\subsection*{The chapter's roadmap}

\Cref{sec:binomial-estimator} retires item~(i): four short
propositions establish that $\hat{\fprate}$ is unbiased, has the
stated variance, equals the MLE, and achieves the Cramer-Rao lower
bound. \Cref{sec:confidence-intervals} retires items~(ii) and~(iii):
recap the Wald form, diagnose its sawtooth, define the Wilson,
Clopper-Pearson, and Jeffreys intervals, and compare their coverage on
the small-$n_0$ regime where Wald fails. \Cref{sec:derived-measures}
retires item~(iv): defines PPV, NPV, recall, and F1 as functions of
$(\fprate, \fnrate, \pi)$, applies the delta method, and works the
prevalence-dependence of PPV as the central caveat.
\Cref{sec:bloom-diagnostics} closes by running the full toolkit on a
deployed Bloom filter with a measured workload, producing a CI for the
filter's PPV and answering a diagnostic question (reduce the
false-positive rate, or accept the current prevalence?).

\subsection*{Notation refresh}

The classification rates introduced in \Cref{sec:quantifying-noise}
and used throughout chapters~1 to~6 are written via the standard
\texttt{alex.sty} macros:
\begin{description}
  \item[$\fprate$ (\texttt{\textbackslash fprate})] false-positive
    rate. The chapter-1 letter was $\varepsilon$.
  \item[$\fnrate$ (\texttt{\textbackslash fnrate})] false-negative
    rate, written $\omega$.
  \item[$\tprate$ (\texttt{\textbackslash tprate})] true-positive rate
    $= 1 - \fnrate$, written $\tau$. Also called \emph{recall} or
    \emph{sensitivity}.
  \item[$\tnrate$ (\texttt{\textbackslash tnrate})] true-negative rate
    $= 1 - \fprate$, written $\eta$. Also called \emph{specificity}.
\end{description}
\noindent The hatted forms $\hat{\fprate}, \hat{\fnrate}, \hat{\tprate},
\hat{\tnrate}$ are the empirical estimators of those rates. Sample
sizes are written $n_0$ for the count of trials with truth $\bot$ and
$n_1$ for the count of trials with truth $\top$, so that
$\hat{\fprate} = \mathrm{FP}/n_0$ and $\hat{\fnrate} = \mathrm{FN}/n_1$.
The Normal-quantile symbol $z_{\alpha/2}$ denotes the upper-$\alpha/2$
critical value of the standard Normal distribution and is used in the
Wald CI of \Cref{sec:sampling-rates} and the Wilson CI of
\Cref{sec:confidence-intervals}. Throughout this chapter we also use
the workload prevalence $\pi = \Pr[X = \top]$ for the fraction of
$\top$-inputs in the inference sample; this is the variable on which
PPV, NPV, and F1 depend and on which $\fprate, \fnrate$ do not.
